\documentclass[10pt,a4paper]{amsart}
\usepackage[margin=19mm]{geometry}
\usepackage{amsmath,amssymb,mathtools}
\usepackage[scr=rsfs,cal=euler]{mathalfa}
\usepackage{xfrac}
\usepackage[unicode,hidelinks,bookmarks=false]{hyperref}
\newcommand{\ie}{i.e.\ }
\newcommand{\cf}{cf.\ }
\newcommand{\resp}{respectively\ }
\newcommand{\Subsection}{\subsection}
\providecommand{\nobreakdash}{\nobreak\hbox{-}}
\setlength{\emergencystretch}{2em}
\multlinegap0pt
\usepackage{amssymb}
\usepackage{tikz}
\usepackage{cancel}
\newtheorem{lemma}{Lemma}
\def\CC{{\mathbb C}}  
\def\Cscr{{\mathscr C}}  
\def\EE{{\mathbb E}} 
\newcommand\Hl{\operatorname{H}}
\newcommand\IH{\operatorname{IH}}
\def\Iscr{{\mathscr I}} 
\def\Lscr{{\mathscr L}}
\def\pt{{\scriptscriptstyle\bullet}}
\def\ssm{\smallsetminus}
\title{Hodge decomposition of $L_2$-cohomology and intersection cohomology of a Shimura variety}
\author{Eduard Looijenga}
\date{Source: arXiv:2501.19253v2 (CC BY 4.0)}
\begin{document}
\maketitle

\noindent\emph{Source and licence.} Hodge decomposition of $L_2$-cohomology and intersection cohomology of a Shimura variety. Version arXiv:2501.19253v2 (CC BY 4.0), distributed by arXiv under the Creative Commons Attribution 4.0 International licence, http://creativecommons.org/licenses/by/4.0/, as recorded in the arXiv licence metadata for this exact version and shown on the abstract page https://arxiv.org/abs/2501.19253v2. Full licence notice: \texttt{licenses/arxiv-2501.19253-CC-BY-4.0.txt}.\par\smallskip

\noindent\textit{Editorial scope.} Counted unit: the article's lemma lemma:limit (source0.tex lines 496--508). The article's complete original proof of it is delivered with it (source0.tex lines 509--515, one \texttt{proof} environment of 7 lines). No mathematics is written, repaired or re-derived: the statement and the proof are the article's own lines, in the article's own notation, and no retained line is substituted or re-flowed.  No further source-local context is needed: every cross-reference of every retained line resolves inside the two delivered spans.  Nothing was omitted from any retained span, so the package carries no omission record.  No retained line contains a citation, so the wrapper carries no bibliography and no bibliography entry.  No retained line contains a figure environment, an image-inclusion command or drawing code, so no image is delivered and the package declares no asset dependency.  Editorial formatting only, in addition: an AMS \texttt{amsart} preamble that restates the article's own declarations and macros used by the retained lines, namely the environments \texttt{lemma} on separate counters, together with the standard packages those lines need.  The retained environments print in this document as Lemma 1 in order of appearance, whereas the article prints the same items with the numbering of its own full document.  The complete native source of the article as distributed in the arXiv e-print for this exact version is bundled unchanged as \texttt{sources/172528/source0.tex}.
\begin{lemma}\label{lemma:limit}  
Assume that we are given $s\in X^*$, an open  neighborhood $U$ of $s$ in $X^*$ and a continuous  proper function 
$d: U\ssm \{s\}\to (0,\infty)$  which exhibits $(U, X\cap U)$ as an open cone with vertex $s$: $d$ extends continuously to $U$ by  mapping $s$ to $\infty$ and  $d: (U\ssm \{s\}, X\cap (U\ssm \{s\}))\to (0,\infty)$ is locally trivial with the local trivializations being on
$ X\cap (U\ssm \{s\})$ a quasi-isometry. 

If the Hodge filtered complex $\Lscr^\pt_{X^*,(2)}(\EE_\CC)$ represents $\Iscr^\pt_{X^*}(\EE_\CC)$ on  $U\ssm \{s\}$, 
then  for  $0\le n<m <\infty$  we have   natural linear  isomorphisms 
\begin{multline}\label{eqn:isos}
\Hl^k_{(2)} (d^{-1}(n,m), \EE_\CC)\cong\IH^k (d^{-1}(n,m), \EE_\CC)\cong\\
\cong  \IH^k(U\ssm \{s\}; \EE_\CC)\cong\Hl^k(i_s^*j_{s*}\Iscr\Cscr_{U\ssm\{s\}}^\pt(\EE_\CC)),
\end{multline}
(where  $j_s:U\ssm \{s\}\subset U$) and  the  Hodge filtration on  the left   defines via the above isomorphisms a family of filtrations on the right with the property that  if we first let $m\to \infty$ and then let $n\to \infty$ we get the natural  Hodge filtration on $\Hl^k(i^*_sj_{s*}\Iscr\Cscr_{U\ssm\{s\}}^\pt(\EE_\CC))$ as a limit.
\end{lemma}

\begin{proof}
Recall that  $\Iscr\Cscr_{U\ssm\{s\}}^\pt(\EE_\CC)$ stands for an object in a derived category of filtered complexes. We are given that the filtered complex $\Lscr^\pt_{U\ssm\{s\},(2)}(\EE_\CC)$ represents that object. Since  the latter   is a complex of  fine sheaves, it follows that the higher direct images $R^kd_*\Lscr^\pt_{U\ssm\{s\},(2)}(\EE_\CC)$ ($k>0$) vanish, so that 
the (zeroth) sheaf direct image $d_*\Lscr^\pt_{U\ssm\{s\},(2)}(\EE_\CC)$ is also a fine complex and represents the direct image 
$d_*\Iscr\Cscr_{U\ssm\{s\}}^\pt(\EE_\CC))$ in a  derived category of  filtered complexes on $(0, \infty)$. So the cohomology of this direct image is a graded local system on $(0, \infty)$ which we can  identify with $R^\pt d_*\Iscr\Cscr_{U\ssm\{s\}}^\pt(\EE_\CC))$. 
In particular, we  get the linear isomorphisms \eqref{eqn:isos}  with the first line even as  isomorphisms of filtered vector spaces. 
It remains to observe  that the definition of a  Hodge module is such that the Hodge  filtration on $\IH^k (d^{-1}(n,m), \EE_\CC)$ has the asserted limiting property.
\end{proof}

\end{document}
